\begin{textsample}{POS Dim 1 – 2017\_c5 – Score 16.00 – 2017\_c5/enade\_2017\_eng\_civ\_9cc7e908b9.txt}  \label{ex:f1_pos_051}
CAPÍTULO 5 DISTRIBUIÇÃO DOS CONCEITOS 5.1 \textbf{PANORAMA} NACIONAL DA DISTRIBUIÇÃO DOS CONCEITOS O Conceito Enade é um indicador de qualidade que avalia os cursos por \textbf{intermédio} dos desempenhos dos estudantes a partir dos resultados do Enade.
Os valores possíveis vão de 1 ( pior situação ) a 5 ( melhor situação ).
Os cursos com apenas um ou sem nenhum concluinte participante não obtêm o Conceito Enade, ficando Sem Conceito ( SC ).
A Tabela 5.1 apresenta a quantidade e distribuição de cursos de Engenharia Civil participantes do Enade/2017, por faixa de conceito e Grande Região.
Enfatiza -se, mais uma vez, que a diferença entre os cursos \textbf{tabulados} neste capítulo e no capítulo 2 corresponde aos cursos Sem Conceito ( SC ), ou seja àqueles, em princípio, sem alunos concluintes que houvessem participado da prova ou àqueles que tivessem tão somente um aluno concluinte18.
Observando -se os dados da Tabela 5.1, nota -se que, dos 539 cursos participantes, 204 ( 37,8 \% ) classificaram -se com conceito 2, o valor modal.
Esse foi também o conceito modal nas regiões Nordeste ( 36,4 \% ) e Sudeste ( 44,8 \% ).
Na região Norte, houve uma distribuição \textbf{bimodal} entre os conceitos 2 e 3 ( 34,5 \% para cada um ).
Nas regiões Sul e Centro-Oeste, o conceito modal foi 3, atribuído, respectivamente, a 45,6 \% e a 33,3 \%.
Em nível nacional, o conceito 3 foi o segundo mais frequente ( 33,4 \%, correspondendo a 180 cursos ) e o conceito 4, o terceiro ( 15,0 \%, correspondendo a 81 cursos ).
Houve, ainda, 28 cursos ( 5,2 \% ) que receberam conceito 5 e 44 cursos ( 8,2 \% ) que receberam conceito 1.
Dos 539 cursos de Engenharia Civil, dois ( 0,4 \% ) ficaram Sem Conceito ( SC ). 18 Esses cursos com somente um concluinte não foram considerados no capítulo 2.
A região Norte participou com 29 cursos ou 5,4 \% do total nacional.
Desses, quatro cursos receberam o conceito 1, o que \textbf{equivale} a 13,8 \% do total regional.
Os conceitos 2 e 3 foram atribuídos a dez cursos, cada um ( 34,5 \% distribuição \textbf{bimodal}, como já comentado ).
A quatro cursos ( 13,8 \% ) atribuiu -se conceito 4, e a um curso apenas ( 3,4 \% ), conceito 5.
Nenhum curso ficou Sem Conceito ( SC ) nessa região.
A região Nordeste participou com 88 cursos ou 16,3 \% do total nacional.
Nessa região, 32 cursos ( 36,4 \% em termos regionais ) obtiveram conceito 2, o conceito modal para a região, como já comentado.
O conceito 1 foi atribuído a seis cursos ( 6,8 \% ).
Ainda nessa região, 25 cursos ( 28,4 \% ) foram avaliados com o conceito 3, 19 cursos ( 21,6 \% ), com conceito 4, e seis cursos ( 6,8 \% ), com conceito 5.
Nenhum curso ficou Sem Conceito ( SC ) nessa região.
Dos 277 cursos participantes da região Sudeste ( 51,4 \% do total nacional ), o conceito 2 foi o conceito modal, como já anteriormente assinalado, tendo sido obtido por 124 cursos ( 44,8 \% ).
O conceito 1 foi atribuído a 32 cursos ( 11,6 \% ), e o conceito 3, a 84 cursos ( 30,3 \% ).
Receberam os conceitos 4 e 5, respectivamente, 21 cursos ( 7,6 \% ) e 14 cursos ( 5,1 \% ).
Nessa região, dois cursos ficaram Sem Conceito ( SC ).
Os 103 cursos da região Sul corresponderam a 19,1 \% do total nacional.
Como já \textbf{indicado}, houve predominância do conceito 3, atribuído a 45,6 \% dos cursos da região ( conceito modal ), ou seja, atribuído a 47 dos 103 cursos participantes na região Sul.
O conceito 2 foi atribuído a 27 cursos ( 26,2 \% ), o conceito 4, a 25 cursos ( 24,3 \% ), e o conceito 5, a quatro cursos ( 3,9 \% ).
Nessa região, nenhum curso ficou Sem Conceito ( SC ), assim como nenhum recebeu conceito 1.
Já dos 42 cursos participantes na região Centro-Oeste ( 7,8 \% do total nacional ), 14 receberam conceito 3, o conceito modal ( 33,3 \% ), como já fora destacado.
A dois cursos foi atribuído conceito 1 ( 4,8 \% ), a 11 cursos ( 26,2 \% ), conceito 2.
Os conceitos \textbf{restantes}, 4 e 5, foram atribuídos, respectivamente, a 12 ( 28,6 \% ) cursos e a três ( 7,1 \% ).
Nessa região, nenhum curso ficou Sem Conceito ( SC ).
O Gráfico 5.1 apresenta a distribuição \textbf{cumulativa} do Conceito Enade segundo a Grande Região.
Num gráfico de distribuição \textbf{cumulativa}, em cada valor de \textbf{abscissa} ( eixo x ), o valor apresentado no eixo das \textbf{ordenadas} ( eixo y ) é o valor acumulado da \textbf{variável}.
Por exemplo, a região Sudeste ( linha vermelha ) apresenta um pouco mais de 87 \% dos cursos com conceito 3 ou abaixo deste valor.
A região Centro-Oeste ( linha preta ), por outro lado, apresenta um pouco mais de 64 \% dos cursos com conceito 3 ou abaixo.
Em linhas gerais, em um gráfico de distribuição \textbf{cumulativa}, \textbf{poligonais} mais à esquerda representam Grandes Regiões com pior distribuição, e \textbf{poligonais} mais à direita, Grandes Regiões com melhores conceitos.
A região Sudeste apresenta os cursos com conceitos mais baixos.
Cursos em IES situadas nas regiões Nordeste ( linha verde ) e Centro-Oeste apresentam uma situação semelhante com uma proporção maior de cursos com conceito máximo, comparadas às das outras regiões. 5.2 CONCEITOS POR CATEGORIA ADMINISTRATIVA, POR MODALIDADE DE ENSINO E GRANDE REGIÃO A Tabela 5.2 apresenta a distribuição dos cursos de Engenharia Civil participantes do Enade/2017, por Categoria Administrativa e Modalidade de ensino, de acordo com os conceitos por eles alcançados, segundo as Grandes Regiões brasileiras.
Dos 539 cursos participantes, 128 ( 23,7 \% ) eram \textbf{ministrados} em Instituições \textbf{Públicas}, e 411 ( 76,3 \% ), em Instituições Privadas.
Quanto à Modalidade de ensino, 537 cursos ( 99,6 \% ) eram presenciais, e apenas dois ( 0,4 \% ), a distância.
De acordo com as informações da Tabela 5.2, em termos nacionais, entre Instituições \textbf{Públicas} e Privadas, observa -se que, dos 28 cursos avaliados com conceito 5, 27 eram oferecidos em IES Públicas, e apenas um era oferecido em IES Privadas.
Dos 128 cursos participantes de IES Públicas, o conceito 4 foi o valor modal, atribuído a 58 cursos ( 45,3 \% ).
Dentre os demais cursos de Instituições \textbf{Públicas} participantes, oito cursos ( 6,3 \% ) foram avaliados com conceito 2, 33 cursos ( 25,8 \% ) obtiveram conceito 3, e 27 cursos ( 21,1 \% ) receberam conceito 5, como já mencionado.
Nessa categoria, um curso ( 0,8 \% ) ficou Sem Conceito ( SC ), e um curso ( 0,8 \% ) ficou com conceito 1.
Na rede privada, o conceito modal foi 2, atribuído a 196 cursos, o correspondente a 47,7 \% dos 411 cursos da categoria.
Dentre os demais cursos participantes, 43 ( 10,5 \% ) receberam conceito 1 e 147 cursos ( 35,8 \% ) receberam conceito 3.
Os conceitos 4 e 5 foram atribuídos, respectivamente, a 23 cursos ( 5,6 \% ) e a um curso ( 0,2 \% ).
Nessa Categoria Administrativa, um curso ( 0,2 \% ) ficou Sem Conceito ( SC ).
Dos 537 cursos participantes da modalidade presencial, o conceito 2 foi o valor modal, atribuído a 202 cursos ( 37,6 \% ).
Dentre os demais cursos nessa modalidade de ensino, 44 cursos ( 8,2 \% ) foram avaliados com conceito 1, 180 cursos ( 33,5 \% ) obtiveram conceito 3, 81 cursos ( 15,1 \% ) receberam conceito 4 e 28 cursos ( 5,2 \% ), conceito 5.
Nessa modalidade, dois cursos ficaram Sem Conceito ( SC ).
Dos dois cursos a distância, ambos receberam conceito 2.
Sendo assim, nessa modalidade, nenhum curso ficou Sem Conceito ( SC ) ou recebeu os demais conceitos.
Na análise por região, observa -se que, na região Norte, as Instituições Privadas participaram com 19 cursos ( 65,5 \% do total regional ), dos quais nenhum ficou Sem Conceito ( SC ), assim como nenhum recebeu conceitos 4 ou 5.
O conceito modal para as Instituições Privadas na região foi 2, atribuído a dez cursos, correspondendo a 52,6 \% dessa categoria na região.
Os conceitos 1 e 3 foram atribuídos, respectivamente, a quatro e a cinco cursos.
As Instituições Públicas participaram com dez cursos na região Norte ( 34,5 \% do total regional ), dos quais cinco obtiveram conceito 3, o valor modal.
Quatro cursos obtiveram conceito 4, e apenas um curso alcançou conceito 5.
Nenhum curso ficou Sem Conceito ( SC ), assim como nenhum curso obteve os conceitos 1 ou 2 nessa combinação de categoria e região.
Todos os 29 cursos oferecidos na região Norte eram presenciais.
Desses, nenhum ficou Sem Conceito ( SC ).
Dez cursos obtiveram conceito 2, outros dez obtiveram conceito 3, valor \textbf{bimodal}.
Quatro cursos obtiveram conceito 1.
O conceito 4 foi atribuído a mais quatro cursos e apenas um curso alcançou o conceito máximo, 5.
Na região Nordeste, a rede privada concentrou 52 dos 88 cursos participantes, o equivalente a 59,1 \% do total da região.
Dentre os cursos oferecidos por IES Privadas no Nordeste, 30 cursos foram avaliados com conceito 2, o conceito modal.
Cinco cursos obtiveram conceito 1, 11 cursos receberam conceito 3, e seis cursos receberam conceito 4.
Nessa combinação de categoria e região, nenhum curso ficou Sem Conceito ( SC ), assim como nenhum alcançou o conceito máximo, 5.
As Instituições \textbf{Públicas} da região Nordeste participaram com 36 cursos ( 40,9 \% do total da região ).
Desses, 14 obtiveram conceito 3, o valor modal, seguidos por outros 13 cursos, que receberam conceito 4.
O conceito 1 foi atribuído a um curso e o conceito 2 a dois cursos.
Seis cursos alcançaram o conceito máximo, 5.
Nenhum curso ficou Sem Conceito ( SC ) nessa região.
Todos os cursos na região foram presenciais : 88.
Para esses, a moda foi o conceito 2, atribuído a 32 cursos.
Os demais cursos presenciais da região foram alocados aos conceitos 1 ( seis cursos ), 3 ( 25 cursos ), 4 ( 19 cursos ) e 5 ( seis cursos ).
Nenhum curso presencial ficou Sem Conceito ( SC ).
Na região Sudeste, a proporção de cursos da rede privada, 87,0 \%, foi mais elevada do que a registrada nas demais regiões brasileiras, correspondendo a 241 dos 277 cursos participantes.
Nessa categoria e região, o conceito modal foi 2, atribuído a 120 cursos.
Um curso ficou Sem Conceito ( SC ).
Os demais foram avaliados com conceito 1 ( 32 cursos ), conceito 3 ( 79 cursos ), conceito 4 ( oito cursos ) e conceito 5 ( um curso ).
Entre os 36 cursos oferecidos por Instituições \textbf{Públicas} na região Sudeste, houve uma distribuição \textbf{bimodal} entre os 4 e 5, atribuídos a 13 cursos, cada um.
Os demais cursos receberam os conceitos 2 ( quatro cursos ) e 3 ( cinco cursos ).
Na região Sudeste, nessa categoria, um curso ficou Sem Conceito ( SC ) e a nenhum curso foi atribuído o conceito 1.
Também na região Sudeste, a maioria dos cursos foram presenciais : 275.
Para esses, a moda foi o conceito 2 com 122 cursos.
Os demais cursos presenciais da região foram alocados aos conceitos 1 ( 32 cursos ), 3 ( 84 cursos ), 4 ( 21 cursos ) e 5 ( 14 cursos ).
Dois cursos presenciais ficaram Sem Conceito ( SC ).
Os dois únicos cursos na modalidade a distância, receberam conceito 2.
Desse modo, nenhum curso a distância da região ficou Sem Conceito ( SC ) ou recebeu os demais conceitos.
As Instituições Privadas concentraram 75 dos 103 cursos participantes da região Sul, 72,8 \% do total regional.
Desses, 41 obtiveram conceito 3, o conceito modal.
Nessa combinação de Categoria Administrativa e Grande Região, os conceitos 2 e 4 foram atribuídos a 26 e a oito cursos, respectivamente.
Nenhum curso recebeu conceitos 1 ou 5.
Nenhum ficou Sem Conceito ( SC ).
As Instituições \textbf{Públicas} na região Sul participaram com 28 cursos, 27,2 \% dos cursos da região.
O conceito modal foi 4, atribuído a 17 cursos.
Aos demais cursos foram atribuídos os conceitos 2 ( um curso ), 3 ( seis cursos ) e 5 ( quatro cursos ).
Nenhum curso da região Sul oferecido por IES Pública ficou Sem Conceito ( SC ) ou recebeu conceito 1.
Também na região Sul, todos os cursos foram presenciais, correspondendo a 103.
Para os cursos presenciais, a moda foi o conceito 3 com 47 cursos.
Os demais cursos presenciais da região foram alocados aos conceitos 2 ( 27 cursos ), 4 ( 25 cursos ) e 5 ( quatro cursos ).
Nenhum curso presencial foi alocado ao conceito 1, assim como nenhum ficou Sem Conceito ( SC ).
Na região Centro-Oeste, 24 dos 42 cursos participantes eram de Instituições Privadas ( 57,1 \% em termos regionais ).
Desses, 11 receberam conceito 3, o conceito modal.
Entre os demais, dois receberam conceito 1, dez receberam conceito 2, e um recebeu conceito 4.
Nenhum curso ficou Sem Conceito ( SC ) ou com conceito 5.
Quanto aos 18 cursos oferecidos por Instituições \textbf{Públicas} na região Centro-Oeste ( 42,9 \% do total regional ), 11 cursos apresentaram conceito 4, conceito modal, e um curso obteve conceito 2.
Os conceitos 3 e 5 foram atribuídos a três cursos, cada um.
Nenhum curso ficou Sem Conceito ( SC ) ou com conceito 1.
Nessa região, todos os cursos foram presenciais e o conceito modal foi 3, atribuído a 14 cursos.
Os demais cursos foram atribuídos aos conceitos 1 ( dois cursos ), 2 ( 11 cursos ), 4 ( 12 cursos ) e 5 ( três cursos ).
Nenhum curso presencial ficou Sem Conceito ( SC ).
O Gráfico 5.2 apresenta a distribuição \textbf{cumulativa} do Conceito Enade segundo a Categoria Administrativa da IES.
Os cursos em IES Públicas ( linha azul ) apresentam uma \textbf{poligonal} mais à direita do que os em IES Privadas ( linha verde ).
Estas apresentam uma porcentagem acumulada de 94,1 \% dos cursos com conceito 3 ou abaixo, enquanto aquelas apresentam apenas 33,1 \% dos cursos nos mesmos conceitos.
Consequentemente, os cursos de IES Públicas apresentam uma distribuição de Conceitos Enade com valores maiores. 5.3 CONCEITOS POR ORGANIZAÇÃO ACADÊMICA E GRANDE REGIÃO Na Tabela 5.3, encontra -se a distribuição dos conceitos atribuídos aos cursos participantes do Enade/2017 na Área de Engenharia Civil, por Organização Acadêmica, segundo as Grandes Regiões brasileiras.
Dos 539 cursos de Engenharia Civil participantes, 248 eram oferecidos em Universidades, 119 em Centros Universitários, 153 em Faculdades e 19 em CEFET/IFET.
Essa distribuição corresponde a, respectivamente, 46,0 \%, 22,1 \%, 28,4 \% e 3,5 \% dos cursos.
De acordo com os dados apresentados, de todos os 28 cursos avaliados com conceito 5, 22 eram vinculados a Universidades.
Os cursos vinculados a esse tipo de Organização Acadêmica tiveram o conceito 3 como conceito modal, atribuído a 85 cursos ( 34,3 \% ).
Os demais cursos vinculados a Universidades receberam os conceitos 1 ( 14 cursos ), 2 ( 73 cursos ), 4 ( 53 cursos ) e 5 ( 22 cursos, como já mencionado ).
Ficou Sem Conceito ( SC ) apenas um curso.
Entre os cursos em Centros Universitários, o conceito modal foi também 3, atribuído a 54 cursos ( 45,4 \% ).
Nenhum curso vinculado a esse tipo de Organização Acadêmica recebeu conceito 5.
Os outros cursos receberam os conceitos 1 ( sete cursos ), 2 ( 47 cursos ) e 4 ( dez cursos ).
Um curso ficou Sem Conceito ( SC ).
Dos 153 cursos mantidos por Faculdades, 83 ( 54,2 \% ) receberam conceito 2, o conceito modal.
O segundo conceito mais frequentemente atribuído aos cursos oferecidos por essa categoria de IES foi o conceito 3, recebido por 37 cursos e 23, conceito 1.
Na sequência, sete cursos obtiveram conceito 4, e três cursos, conceito 5.
Nenhum curso ficou Sem Conceito ( SC ).
Dos 19 cursos mantidos por CEFET/IFET, 11 ( 57,9 \% ) receberam conceito 4, o conceito modal.
O segundo conceito mais frequentemente atribuído aos cursos oferecidos por essa categoria de IES foi o conceito 3, recebido por quatro cursos.
Três cursos receberam conceito 5.
O conceito 2 foi atribuído a um curso.
Nenhum curso ficou Sem Conceito ( SC ) ou recebeu conceito 1.
Considerando -se separadamente as regiões brasileiras, verifica -se que, na região Norte, as Universidades participaram com 11 dos 29 cursos participantes.
Cinco cursos receberam conceito 3, conceito modal.
Os demais cursos receberam os conceitos 1 ( dois cursos ), 2 ( um curso ), 4 ( dois cursos ) e 5 ( um curso ).
Nenhum curso ficou Sem Conceito ( SC ) nessa combinação de categoria e região.
Os Centros Universitários da região Norte foram representados por cinco cursos, aos quais foram atribuídos os conceitos 2 ( valor modal, recebido por quatro cursos ) e 3 ( recebido por um curso ).
Nenhum curso ficou Sem Conceito ( SC ) ou recebeu os demais conceitos nessa combinação de categoria e região.
As Faculdades participaram com 11 cursos na região Norte que receberam os conceitos 1 ( dois cursos ), 2 ( cinco cursos, conceito modal ) e 3 ( quatro cursos ).
Nenhum curso ficou Sem Conceito ( SC ) e nenhuma Faculdade da região obteve conceito 4 e 5.
Os CEFET/IFET participaram com dois cursos, que receberam conceito 4.
Na região Nordeste, as Universidades participaram com 41 dos 88 cursos da Área de Engenharia Civil oferecidos na região, dos quais nenhum ficou Sem Conceito ( SC ).
Houve uma distribuição modal entre os conceitos contíguos 3 e 4, cada um representado por 13 cursos.
Apenas um curso obteve conceito 1, nove obtiveram conceito 2, e cinco, conceito máximo 5.
Os Centros Universitários \textbf{contaram} com 14 cursos participantes na região Nordeste.
Esses obtiveram conceito 1 ( um curso ), 2 ( conceito modal, obtido por seis cursos ), 3 ( cinco cursos ), e 4 ( dois cursos ).
Nenhum curso oferecido por Centros Universitários no Nordeste ficou Sem Conceito ( SC ) ou recebeu conceito 5.
As Faculdades foram representadas por 28 cursos na região Nordeste.
O conceito 1 foi atribuído a quatro cursos, o conceito 2, a 17 cursos, o valor modal, o conceito 3 a quatro cursos, e o conceito 4 a três cursos.
Na região Nordeste, nenhum curso oferecido por Faculdades recebeu conceito 5, tampouco ficou Sem Conceito ( SC ).
Os CEFET/IFET participaram com cinco cursos, os quais receberam o conceito modal 3 ( três cursos ).
Os conceitos 4 e 5 foram atribuídos a um curso, cada um.
Nenhum curso ficou Sem Conceito ( SC ), ou recebeu os conceitos 1 ou 2.
Na região Sudeste, as Universidades concentraram 119 dos 277 cursos de Engenharia Civil da região.
Entre os cursos em Universidades, nessa região, o conceito modal foi 2, atribuído a 48 cursos.
Os demais cursos receberam os conceitos 1 ( 11 cursos ) e 3 ( 37 cursos ).
Os conceitos 4 e 5 foram atribuídos a 11 cursos, cada um.
Apenas um curso ficou Sem Conceito ( SC ) e nenhum curso recebeu o conceito máximo, 5.
Os Centros Universitários participaram com 69 cursos na região Sudeste, dos quais 30 obtiveram o conceito modal 3.
O conceito 1 foi atribuído a seis cursos, o conceito 2 a 29 cursos, e o conceito 4, a três cursos.
Nenhum curso alcançou conceito 5.
Um curso apenas ficou Sem Conceito ( SC ) nessa combinação de organização acadêmica e região.
Já as Faculdades foram representadas por 85 cursos na região Sudeste, dos quais nenhum ficou Sem Conceito ( SC ).
O conceito modal foi o conceito 2, atribuído a 47 cursos.
O conceito 1 foi atribuído a 15 cursos, e o conceito 3, a 17 cursos.
Já os conceitos 4 e 5 foram atribuídos a três cursos, cada um.
Os CEFET/IFET participaram com quatro cursos, os quais receberam conceito 4.
Sendo assim, nenhum curso ficou Sem Conceito ( SC ), ou recebeu os demais conceitos.
Dos 103 cursos da região Sul, 59 eram de Universidades, para os quais o conceito modal foi 3, atribuído a 25 cursos.
Nesse tipo de organização, nenhum curso ficou Sem Conceito ( SC ) ou recebeu conceito 1.
O conceito 2 foi atribuído a 11 cursos, o conceito 4, a 20 cursos, e o conceito 5, a três cursos.
Dos 20 cursos participantes de Centros Universitários da região Sul, nenhum ficou Sem Conceito ( SC ) ou recebeu conceitos 1 ou 5.
O conceito modal, 3, foi atribuído a 12 cursos.
Os conceitos \textbf{restantes}, 2 e 4, foram atribuídos a quatro cursos, cada um.
Quanto aos 22 cursos vinculados a Faculdades na região Sul, 11 receberam o conceito modal 2.
O conceito 3 foi atribuído a dez cursos, e o conceito 4, a um curso.
Nenhum curso ficou Sem Conceito ( SC ) ou recebeu conceito 1 ou 5.
Nessa região, foram avaliados dois cursos em CEFET/IFET, dos quais um recebeu conceito 2, e outro recebeu conceito 5.
Na região Centro-Oeste, 18 dos 42 cursos eram de Universidades.
Nesse tipo de organização, nenhum curso ficou Sem Conceito ( SC ) ou recebeu conceito 1.
O conceito 2 foi atribuído quatro cursos, o conceito 3, a cinco cursos.
O conceito 4, conceito modal, foi atribuído a sete cursos.
Além disso, dois cursos alcançaram o conceito máximo, 5.
Os Centros Universitários da região Centro-Oeste \textbf{contaram} com 11 cursos, que obtiveram os conceitos 3 ( seis cursos, conceito modal ), 2 ( quatro cursos ) e 4 ( um curso ).
Nenhum curso ficou Sem Conceito ( SC ) ou recebeu conceito 1 e 5.
Dos sete cursos oferecidos em Faculdades na região Centro-Oeste, três receberam o conceito modal 2.
Os conceitos 1 e 3 foram atribuídos a dois cursos, cada um.
Nenhum curso ficou Sem Conceito ( SC ), assim como nenhum curso recebeu os conceitos 4 e 5.
Nessa região, foram avaliados seis cursos em CEFET/IFET, dos quais quatro receberam o conceito modal 4.
Os conceitos 3 e 5 foram atribuídos a um curso, cada um.
Os demais conceitos não foram atribuídos a nenhum curso na região e nenhum curso ficou Sem Conceito ( SC ).
O Gráfico 5.3 apresenta a distribuição \textbf{cumulativa} do Conceito Enade segundo Organização Acadêmica da IES.
Os cursos em CEFET/IFET ( linha \textbf{roxa} ) apresentam uma proporção \textbf{cumulativa} de 26,3 \% no conceito 3 ou abaixo, ou seja, uma \textbf{poligonal} mais à direita do que as demais, \textbf{denotando} uma melhor distribuição de Conceitos Enade.
Em IES de Centros Universitários ( linha verde ) todos os cursos têm conceitos 4 ou abaixo disso, já que a \textbf{poligonal} \textbf{cumulativa} alcança 100 \% para o conceito 4.
IES de Faculdades ( linha vermelha ) apresenta uma maior proporção de cursos com conceitos mais baixos ( 69,3 \% no conceito 2 ou abaixo ) e em Universidades ( linha azul ) apresentam uma situação intermediária.

% matched lemmas: abscissa, bimodal, contar, cumulativo, denotar, equivaler, indicado, intermédio, ministrar, ordenado, panorama, poligonal, pública, restante, roxa, tabular, variável
\end{textsample}
